\section*{Capítulo \ref{ch:b2:reaction-free-energy} --- Entropia e energia de Gibbs de reação}

\begin{solution}{exo:b2:reaction-free-energy:1}
(a) Três mols de gás passam a líquido: $\Delta_r S^\circ < 0$. (b) Forma-se
um gás: $> 0$. (c) Uma molécula de gás dá duas: $> 0$. (d) Íons em solução
formam um cristal: $< 0$. (e) Fusão: $> 0$. (f) $\Delta\nu_{\text{gás}} = 0$:
pequena e, aqui, positiva (\qty{+20}{J/(K.mol)} pelas tabelas, sendo \ce{HCl}
menos simétrico que \ce{H2} e \ce{Cl2}).
% ledger: s0:HCl_g, s0:H2_g, s0:Cl2_g
\end{solution}

\begin{solution}{exo:b2:reaction-free-energy:2}
$\Delta_r S^\circ = 2(192.77) - 191.61 - 3(130.68) = \qty{-198.1}{J/(K.mol)}$;
$\Delta_r G^\circ = -91.9 - 298.15 \times (-0.1981) = \qty{-32.8}{kJ/mol}$:
\omterm{def:b2:reaction-free-energy:exergonic}{exergônica} a \qty{298}{K}, embora o termo entrópico (\qty{+59.1}{kJ/mol})
jogue contra.
\end{solution}

\begin{solution}{exo:b2:reaction-free-energy:3}
$\Delta_r S^\circ = 69.95 - 130.68 - \tfrac12(205.15) = \qty{-163.3}{J/(K.mol)}$;
$\Delta_r G^\circ = -285.83 + 298.15 \times 0.1633 = \qty{-237.14}{kJ/mol}$.
É a reação de formação da água líquida (um mol a partir dos elementos em
seus estados de referência); logo, sua \omterm{def:b2:reaction-free-energy:reaction-gibbs}{energia de Gibbs padrão de reação} é,
por definição, $\Delta_f G^\circ(\ce{H2O}, \text{l})$; as tabelas dão o mesmo
valor.
\end{solution}

\begin{solution}{exo:b2:reaction-free-energy:4}
Cerca de \qty{42}{J/(K.mol)} a \qty{298}{K} e \qty{87}{J/(K.mol)} a
\qty{1000}{K} (líquido). Saltos: \qty{10.6}{J/(K.mol)} a \qty{692.7}{K} e
\qty{97.7}{J/(K.mol)} a \qty{1180.2}{K}. $\Delta_{\text{fus}}H = 10.6 \times
692.7 = \qty{7.3}{kJ/mol}$ and $\Delta_{\text{vap}}H = 97.7 \times 1180.2 =
\qty{115}{kJ/mol}$.
\end{solution}

\begin{solution}{exo:b2:reaction-free-energy:5}
$\Delta_r H^\circ = \qty{44.00}{kJ/mol}$, $\Delta_r S^\circ = 188.84 - 69.95 =
\qty{118.9}{J/(K.mol)}$, $\Delta_r G^\circ = 44.00 - 298.15 \times 0.1189 =
\qty{+8.55}{kJ/mol}$. A \qty{298}{K}, a transformação não é reversível (a água
líquida e o vapor a \qty{1}{bar} não estão em equilíbrio); logo $\Delta S \ne
Q/T$: $\Delta_r H^\circ/T = \qty{147.6}{J/(K.mol)}$. O $\Delta_r G^\circ$
positivo significa que, a \qty{298}{K}, a água não evapora em uma atmosfera
de vapor puro a \qty{1}{bar}: o vapor se condensa. Os dois
estão em equilíbrio onde $\Delta_r G^\circ = 0$: $T = 44\,004/118.9 =
\qty{370}{K}$, a \qty{3}{K} do ponto de ebulição normal
(\qty{373}{K}): a \omterm{def:b2:reaction-free-energy:ellingham-approximation}{aproximação de Ellingham} é boa em um intervalo de \qty{75}{K}.
\end{solution}

\begin{solution}{exo:b2:reaction-free-energy:6}
$T_i = 57\,100/175.7 = \qty{325}{K}$. Abaixo dela, $\Delta_r G^\circ > 0$ e o
\ce{N2O4}, incolor, é favorecido; resfriar o tubo desloca o gás no sentido do
\ce{N2O4}, e a cor castanha do \ce{NO2} desbota (a relação quantitativa com
a composição está no \cref{ch:b2:equilibrium-shifts}).
\end{solution}

\begin{solution}{exo:b2:reaction-free-energy:7}
$\Delta_r S^\circ = -\dd\Delta_r G^\circ/\dd T = -(-4.0 + 12.0)/100 =
\qty{-0.080}{kJ/(K.mol)} = \qty{-80}{J/(K.mol)}$ and $\Delta_r H^\circ =
\Delta_r G^\circ + T\Delta_r S^\circ = -12.0 + 300 \times (-0.080) =
\qty{-36.0}{kJ/mol}$. Verificação: $\Delta_r G^\circ/T$ vai de $-0.0400$ a
$\qty{-0.0100}{kJ/(K.mol)}$, inclinação $3.00 \times 10^{-4}$ por kelvin, igual a
$-\Delta_r H^\circ/T^2 = 36.0/346^2 = 3.00 \times 10^{-4}$ na temperatura
média geométrica $\sqrt{300 \times 400} = \qty{346}{K}$.
\end{solution}

\begin{solution}{exo:b2:reaction-free-energy:8}
$\Delta_r S^\circ = 186.25 - 5.74 - 2(130.68) = \qty{-80.8}{J/(K.mol)}$,
$\Delta_r G^\circ = -74.87 + 298.15 \times 0.0808 = \qty{-50.8}{kJ/mol}$,
o $\Delta_f G^\circ$ tabelado do metano. Como os dois termos têm o mesmo
sinal, o metano se torna instável acima de $T_i = 74\,870/80.8 \approx
\qty{926}{K}$ (\omterm{def:b2:reaction-free-energy:ellingham-approximation}{aproximação de Ellingham}): o metano quente se decompõe em
carbono e hidrogênio, uma via para obter ambos.
\end{solution}

\begin{solution}{exo:b2:reaction-free-energy:9}
Substituição: \ce{CH3CHBrCH3 + OH- -> CH3CH(OH)CH3 + Br-}; eliminação:
\ce{CH3CHBrCH3 + OH- -> CH2=CHCH3 + H2O + Br-}. A diferença entre as duas
\omterm{def:b2:reaction-free-energy:reaction-gibbs}{energias de Gibbs padrão de reação} é $\Delta(\Delta_r G^\circ) = 40 - T \times
0.090$, negativa acima de $40\,000/90 \approx \qty{440}{K}$. A eliminação forma
uma molécula a mais que a substituição: sua entropia é maior, e o termo
entrópico cresce proporcionalmente a $T$.
\end{solution}

\begin{solution}{exo:b2:reaction-free-energy:10}
Cada uma das $N_A$ moléculas tem 2 orientações: $W = 2^{N_A}$, logo $S =
k_B\ln 2^{N_A} = R\ln 2 = \qty{5.76}{J/(K.mol)}$. O cristal não é
perfeitamente ordenado: a baixa temperatura, as moléculas ficam congeladas
em suas orientações aleatórias e não conseguem atingir o arranjo ordenado
único que a terceira lei supõe.
\end{solution}

\begin{solution}{exo:b2:reaction-free-energy:11}
Integre $\dd\Delta_r S^\circ/\dd T = c/T$: $\Delta_r S^\circ(T) = \Delta_r
S^\circ(T_0) + c\ln(T/T_0)$; Kirchhoff dá $\Delta_r H^\circ(T) = \Delta_r
H^\circ(T_0) + c(T - T_0)$; subtraia $T\Delta_r S^\circ(T)$. Para o calcário a
\qty{1100}{K}: Ellingham $178.32 - 1100 \times 0.16064 = \qty{1.62}{kJ/mol}$;
correção $c[T - T_0 - T\ln(T/T_0)] = -0.00195 \times (801.85 - 1436.3) =
\qty{+1.24}{kJ/mol}$: $\Delta_r G^\circ = \qty{2.86}{kJ/mol}$. O erro,
\qty{1.2}{kJ/mol}, é menos de um por cento de $\Delta_r H^\circ$, mas, perto
da \omterm{def:b2:reaction-free-energy:inversion-temperature}{temperatura de inversão}, ele decide o sinal.
\end{solution}

\begin{solution}{exo:b2:reaction-free-energy:12}
(a) $\Delta_r H^\circ = -763.2 + 944.0 = \qty{+180.8}{kJ/mol}$, $\Delta_r
S^\circ = 353.2 + 205.15 - 50.62 - 2(223.08) = \qty{+61.6}{J/(K.mol)}$;
$\Delta_r G^\circ(1000) = 180.8 - 61.6 = \qty{+119.2}{kJ/mol}$: \omterm{def:b2:reaction-free-energy:exergonic}{endergônica},
não há cloração. (b) Para \ce{C + O2 -> CO2}, $\Delta_r G^\circ(1000) =
-393.5 - 2.9 = \qty{-396.4}{kJ/mol}$. A soma, \ce{TiO2(s) + 2Cl2(g) + C(s)
-> TiCl4(g) + CO2(g)}, tem $\Delta_r G^\circ = \qty{-277.2}{kJ/mol}$: o carbono
consome o dioxigênio formado pela primeira reação, e a reação acoplada é
fortemente \omterm{def:b2:reaction-free-energy:exergonic}{exergônica}.
\end{solution}

\begin{solution}{pb:b2:reaction-free-energy:1}
\textbf{1.} \ce{CaCO3(s) -> CaO(s) + CO2(g)}.
\textbf{2.} $\Delta_r H^\circ = -635.09 - 393.51 + 1206.92 =
\qty{+178.3}{kJ/mol}$: \omterm{def:b2:reaction-enthalpy:exothermic}{endotérmica}.
\textbf{3.} $\Delta_r S^\circ = 39.75 + 213.79 - 92.9 = \qty{+160.6}{J/(K.mol)}$,
positiva porque um gás se forma a partir de um sólido.
\textbf{4.} $\Delta_r G^\circ = 178.32 - 298.15 \times 0.16064 =
\qty{+130.4}{kJ/mol}$.
\textbf{5.} Sim: é preciso $\Delta_r G < 0$ para a decomposição, e aqui ele
é fortemente positivo; o dióxido de carbono do ar, muito abaixo de \qty{1}{bar},
diminui $\Delta_r G$, mas muito menos que \qty{130}{kJ/mol} (próximo capítulo).
\textbf{6.} $\Delta_r C^\circ_p = 42.80 + 37.13 - 81.88 =
\qty{-1.95}{J/(K.mol)}$, minúsculo: $\Delta_r H^\circ$ e $\Delta_r S^\circ$
quase não variam com $T$.
\textbf{7.} $\Delta_r G^\circ(T) = 178.32 - 0.16064\,T$ (\unit{kJ/mol}).
\textbf{8.} \qty{+17.7}{kJ/mol} a \qty{1000}{K}; \qty{-30.5}{kJ/mol} a
\qty{1300}{K}.
\textbf{9.} $T_i = 178\,320/160.64 = \qty{1110}{K}$.
\textbf{10.} $\Delta_r G^\circ/T = 178.32/T - 0.16064$, cuja derivada é
$-178.32/T^2 = -\Delta_r H^\circ/T^2$.
\textbf{11.} Em $T_i$, $c[T_i - T_0 - T_i\ln(T_i/T_0)] = -0.00195 \times
(811.9 - 1459.4) = \qty{+1.26}{kJ/mol}$; $\Delta_r G^\circ$ se anula
$1260/160.6 \approx \qty{8}{K}$ mais acima, em cerca de \qty{1118}{K}: menos
de um por cento.
\textbf{12.} A \qty{1000}{K}, $\Delta_r G > 0$: a decomposição não pode
ocorrer (a cal absorveria dióxido de carbono). A \qty{1300}{K}, $\Delta_r G <
0$: o calcário se decompõe.
\textbf{13.} $\Delta_r G$ se torna positivo: a cal capta o dióxido de carbono
e volta a ser carbonato.
\textbf{14.} $\delta S_{\text{c}} = -\Delta_r G\,\dd\xi/T = 30\,506/1300 =
\qty{23.5}{J/K}$ por mol.
\textbf{15.} Sim: o \cref{ch:b2:chemical-potential} mostrará que
$\Delta_r G = \Delta_r G^\circ + RT\ln(p_{\ce{CO2}}/p^\circ)$, que diminui
quando a pressão de dióxido de carbono está abaixo de \qty{1}{bar}, de modo
que a decomposição começa abaixo de $T_i$.
\textbf{16.} $n = 10^6/56.08 = \qty{1.783e4}{mol}$.
\textbf{17.} $1.783 \times 10^4 \times 178.32 = \qty{3.18e6}{kJ}$, isto é,
\qty{3.18}{GJ} por tonelada.
\textbf{18.} $1.783 \times 10^4 \times 44.01 = \qty{785}{kg}$ de \ce{CO2}.
\textbf{19.} Calor a fornecer: $3.18 \times 10^6/0.50 = \qty{6.36e6}{kJ}$, logo
$6.36 \times 10^6/802.3 = \qty{7.93e3}{mol}$ de metano, \qty{127}{kg}.
\textbf{20.} $7.93 \times 10^3 \times 44.01 = \qty{349}{kg}$ de \ce{CO2}.
\textbf{21.} $785 + 349 = \qty{1134}{kg}$ por tonelada de cal, dos quais 69\,\%
vêm do calcário.
\textbf{22.} O dióxido de carbono do calcário é fixado pela estequiometria,
seja qual for a fonte de calor do forno; só a parte do combustível pode ser
reduzida.
\textbf{23.} O calcário se decompõe sob \qty{1}{bar} de dióxido de carbono
acima de sua \omterm{def:b2:reaction-free-energy:inversion-temperature}{temperatura de inversão}, $\boldsymbol{T_i \approx \qty{1.11e3}{K}}$
(cerca de \qty{840}{\celsius}).
\end{solution}
